\documentclass[11pt]{article}

\usepackage{amsmath, amssymb, amsthm}
\usepackage{fullpage}
\usepackage{hyperref}


\usepackage{mathtools}

% Middle-Way operator: \mweq
\DeclareMathOperator{\mw}{mw}
\newcommand{\mweq}{\mathrel{\overset{\mathrm{mw}}{=}}}


\title{Physical Constants as Middle-Way Fixed Points:\\
A Dynamical Resolution of the Fundamental Constant Problem}

\author{Masaomi Chiba}

\date{2025-12-10}

\begin{document}

\maketitle

\begin{abstract}
Modern physics contains a number of unexplained empirical constants,
including the speed of light $c$, the Planck constant $h$, the gravitational
constant $G$, the cosmological constant $\Lambda$, and various coupling
constants. These constants are treated as fixed, universal quantities,
yet no fully satisfactory explanation has been given for why they take
their observed values.

This paper proposes a new approach: physical constants are not static,
primitive entities, but are instead \emph{Middle-Way Fixed Points}
that arise from the dynamical interaction of uncertainty (emptiness),
stochastic proposals (ranki), and structural constraints (dependent
origination). These components form a general coherence functional
whose minima define the apparent ``constants.'' The resulting framework
resolves the Fundamental Constant Problem by replacing static equality
with a dynamic equivalence operator, written $\mweq$, which captures
convergence to minimal-coherence fixed points rather than strict identity.

This view parallels the resolution of classical logical paradoxes
(Godel incompleteness) and computational barriers (such as
the P vs NP problem in the companion paper ``P $\mweq$ NP'').
\end{abstract}

\section{Introduction}

Why do the constants of nature take the values they do?
The question is older than relativity or quantum theory, and
persists across modern physics. Despite impressive theoretical
developments, no derivation exists for the numerical values of
$c$, $h$, $G$, $\Lambda$, or the masses of elementary particles.

Traditionally, physics assumes a fixed observer frame and a fixed
background structure. Under these assumptions, constants necessarily
appear arbitrary. This mirrors earlier difficulties in logic and
computation, where static frameworks produced undecidability and
incompleteness phenomena.

We propose that the static assumption itself is the root problem.
If the observer is not fixed---that is, if the measurement frame
is part of the dynamical system---then constants are not primitive
but emergent. They arise as stable fixed points in a coherence-minimizing
process.

\section{Dynamic Structure: Emptiness, Ranki, and Dependent Origination}

We model the emergence of physical constants by three structural
components:

\begin{enumerate}
    \item \textbf{Emptiness (E)}: an initial space of uncertainty,
    representing the degrees of freedom in the measurement and physical
    configuration.

    \item \textbf{Ranki (R)}: stochastic proposals generated from $E$,
    analogous to non-deterministic sampling, evolutionary search, or
    Monte Carlo exploration.

    \item \textbf{Dependent Origination (D)}: a constraint functional
    that assigns coherence energy to each proposal.
\end{enumerate}

Let $x$ be a candidate value (e.g., for $c$ or $G$).
Define the coherence energy functional:
\[
E(x) = SC(x) + CL(x) + OD(x),
\]
where $SC$ (self-consistency), $CL$ (cross-linearity),
and $OD$ (observer-dependence) are non-negative terms.

A physical constant is then defined not by a static equation,
but as a \emph{minimal-coherence fixed point}:
\[
x^{\ast} = \arg \min_{x} E(x).
\]

\section{The Middle-Way Equivalence Operator}

The traditional equality sign is inadequate for dynamical systems in
which the observer and the observed mutually constrain each other.
We therefore introduce the Middle-Way equivalence operator:
\[
$$ A $\mweq$ B \quad \text{iff} \quad A \text{ and } B $$
\text{ converge to the same minimal-coherence fixed point}.
\]

This operator generalizes classical equality by grounding equivalence
in dynamical convergence rather than static identity.

Applying $\mweq$ to physical constants means that
``$c$ is constant'' is not a statement of rigid identity, but of
stability under dynamical observation. This resolves longstanding
confusions about whether $c$ ``really'' changes or is an illusion
of measurement: both sides converge in the limit
to the same coherence minimum.

\section{Application to the Speed of Light}

The speed of light has long been a paradoxical quantity.
It is invariant across observers, yet emerges from the electromagnetic
properties of the vacuum:
\[
c = \frac{1}{\sqrt{\mu_{0} \epsilon_{0}}}.
\]

Under our framework, this is natural. $c$ is the minimal-coherence
fixed point of:

\begin{itemize}
    \item the structure of spacetime geometry,
    \item electromagnetic propagation,
    \item information-theoretic constraints on causality,
    \item and the dynamical observer configuration.
\end{itemize}

Thus:
\[
 c \mweq x^{\ast},
\]
and the speed of light appears constant because it is
the stable equilibrium of the coherence functional.

\section{Resolution of the Fundamental Constant Problem}

The proposed framework resolves the Fundamental Constant Problem
without fine-tuning or anthropic reasoning:

\begin{enumerate}
    \item Constants are not primitive, but emergent.
    \item Their values are determined by coherence minimization.
    \item The observer is part of the system, so constants
    must be understood dynamically.
    \item Classical equality is replaced by Middle-Way
    equivalence: $A \mweq B$.
\end{enumerate}

In this view, physical constants are not arbitrary.
They are the stable fixed points of a deep dynamical structure.

\section{Conclusion}

We have proposed that physical constants arise as Middle-Way fixed
points---minimal-coherence equilibria determined by the interaction
of uncertainty, stochastic proposals, and structural constraints.
This approach replaces static equality with a dynamic equivalence
relation and resolves longstanding foundational puzzles.

This framework parallels similar resolutions in logic and
computation, suggesting a unified dynamical architecture for
emergent structure across scientific domains.

\bibliographystyle{plain}
\bibliography{refs}

\end{document}

